\subsection{Incidencias entre mitades de emisiones regionales}
\label{sec:rec27-empalmes}

Sea $\mathcal U_6$ la imagen de las seis ventanas consecutivas del emisor
de \ref{act21:tpk:emision}. Para $u=(u_1,\ldots,u_6)$ definimos
$A(u)=(u_1,u_2,u_3)$ y $B(u)=(u_4,u_5,u_6)$.
El grafo de mitades tiene como vértices esos triples y una arista no
orientada entre $A(u)$ y $B(u)$ por cada emisión.

\begin{proposition}[Ciclo de empalmes regionales]
Los vértices
\[
\begin{aligned}
A&=(830,943,830),&B&=(109,252,252),\\
C&=(810,923,870),&D&=(169,272,272),\\
E&=(850,963,890),&F&=(149,252,212)
\end{aligned}
\]
forman el ciclo $A-B-C-D-E-F-A$.
\end{proposition}
\begin{proof}
Escribimos cada semilla como $(i,j,d)$, con posiciones residuales entre
$0$ y $8$ y direcciones $N=(-1,0)$, $E=(0,1)$ y $S=(1,0)$.
Las siguientes parejas de semillas producen, mediante la recurrencia de
54 instantes, las seis emisiones indicadas:
\[
\begin{array}{c|c|c}
\text{emisión}&s_+&s_\times\\\hline
A\mid B&(0,6,E)&(0,6,N)\\
B\mid C&(0,6,N)&(1,5,E)\\
C\mid D&(0,6,E)&(0,1,S)\\
D\mid E&(0,6,N)&(0,4,S)\\
E\mid F&(0,6,E)&(0,3,E)\\
F\mid A&(0,6,N)&(0,1,S)
\end{array}
\]
Cada fila es un testigo de pertenencia a $\mathcal U_6$, obtenido por
evaluación sucesiva de los dos acumuladores. Al separar cada emisión en
sus dos mitades se obtienen exactamente las seis aristas del ciclo.
\end{proof}

En la clasificación de \eqref{eq:palabras-regionales},
\(A\mid B\), \(C\mid D\) y \(E\mid F\) producen respectivamente
las palabras \(w^{\rm au}\), \(w^{\rm cl}\) y \(w^{\rm pr}\).
Los lectores regionales identifican posteriormente estas salidas con
\(\varphi\), una región de \(\pi\) y \(e\).
Las otras tres emisiones hacen explícito su empalme en el grafo.
El resultado describe incidencias entre salidas del emisor. La realización
temporal de una incidencia utiliza, adicionalmente, la orientación, la fase
y la memoria del estado enriquecido. La incorporación del anclaje APP y
de las cinco regiones de $\pi$ impone condiciones de recorrido adicionales;
el ciclo anterior conserva su función como testigo local de empalme.

El ciclo de seis aristas pertenece a este grafo no dirigido de
semiemisiones. Un recorrido temporal enriquecido conserva además los
datos de anclaje, orientación y cierre de su construcción. La incidencia
entre dos mitades y el transporte de un estado tienen así dominios y
condiciones de composición explícitos.
